\section{Cross-area research projects}
\label{app:projects}
% DRAFT (Claude): awaits the author's rewrite. The table and numbers are written by
% `markov-lean paper` from results/cases/<name>/summary.json, results/pfr/summary.json, each
% project's switches/ run on judge labels, and results/cases/compare.json
% (scripts/after_case_judges.sh after the judges).

Mathlib is a library: its proofs are short and its lemmas are filed by area.
To see how proofs behave where a result is known for joining fields, we
extracted every tactic step of four research formalizations with the same
rules as LeanDojo: sphere packing in dimension 8, the prime number theorem and
its extensions (PNT+), the polynomial Freiman--Ruzsa conjecture (PFR), and, as a
control that stays in one field, Carleson's theorem in harmonic analysis
\todo{cite the four projects}. Table~\ref{tab:projects} gives two kinds of area
label for each. Premise areas are free: a step's area is the majority Mathlib
folder of the lemmas it cites, mapped to \textsc{msc} classes, optionally also
counting the project's own lemmas by the file that defines them. Judge areas are
the two judges' agreed \textsc{msc} codes, as in Section~\ref{sec:switches}.

Premise areas switch about as often in the control
(\resCaseCarlesonPremRate{} switches per move) as in sphere packing
(\resCaseSpherePremRate{}) or PNT+ (\resCasePNTPremRate{}): long research
proofs cite many Mathlib folders whatever their subject, so premise labels
cannot tell a proof that crosses fields from one that only uses a lot of
library. The judge switch rate minus the control's is
\resCaseSphereVsCtlRate{} [\resCaseSphereVsCtlRateLo, \resCaseSphereVsCtlRateHi]
switches per move in sphere packing ($p = \resCaseSphereVsCtlRateP$),
\resCasePNTVsCtlRate{} [\resCasePNTVsCtlRateLo, \resCasePNTVsCtlRateHi] in PNT+
($p = \resCasePNTVsCtlRateP$) and \resCasePFRVsCtlRate{}
[\resCasePFRVsCtlRateLo, \resCasePFRVsCtlRateHi] in PFR
($p = \resCasePFRVsCtlRateP$).

\begin{table*}[t]
  \caption{Area switching in four research formalizations. Premise areas: all
  extracted steps, switches between consecutive labelled steps, pooled over
  moves; ``with own'' also labels the project's own lemmas by their file.
  Judge areas: proofs with at least two judge-labelled steps, each proof
  weighted equally, with the mean under the within-proof shuffle and the ratio
  of the two (below 1 means a proof's steps in one area come in runs). Carleson's
  theorem stays in harmonic analysis and is the control.}
  \label{tab:projects}
  \centering\small
  \setlength{\tabcolsep}{4pt}
  \resultstable{table_projects}
\end{table*}
